\documentclass{article}

% if you need to pass options to natbib, use, e.g.:
%     \PassOptionsToPackage{numbers, compress}{natbib}
% before loading neurips_2024


% ready for submission - DOUBLE BLID
% \usepackage{neurips_2024}


% to compile a preprint version, e.g., for submission to arXiv, add add the
% [preprint] option:
%     \usepackage[preprint]{neurips_2024}


% to compile a camera-ready version, add the [final] option, e.g.:
% \usepackage{neurips_2025}
\usepackage{iclr2026_conference,times}
% \iclrfinalcopy

% to avoid loading the natbib package, add option nonatbib:
%    \usepackage[nonatbib]{neurips_2024}

\usepackage[utf8]{inputenc} % allow utf-8 input
\usepackage[T1]{fontenc}    % use 8-bit T1 fonts
\usepackage{hyperref}       % hyperlinks
\usepackage{url}            % simple URL typesetting
\usepackage{booktabs}       % professional-quality tables
\usepackage{amsfonts}       % blackboard math symbols
\usepackage{nicefrac}       % compact symbols for 1/2, etc.
\usepackage{microtype}      % microtypography
\usepackage{xcolor}         % colors
\usepackage{subfiles}
\usepackage{lipsum}
\usepackage{graphicx} % For check and cross symbols
\usepackage{pifont}
\usepackage{array}
\usepackage{float}
\usepackage{amsmath}
\usepackage{algorithm}
\usepackage{algorithmicx}
\usepackage{algpseudocodex}
\usepackage{enumitem}
\usepackage{titlesec}
\usepackage{booktabs}
\usepackage{multirow}
\usepackage{graphicx}
\usepackage{makecell}
\usepackage{geometry}
\usepackage{listings}
\usepackage{xcolor}

\lstdefinestyle{mystyle}{
    basicstyle=\ttfamily\small,
    breaklines=true,
    frame=single,
    columns=fullflexible,
    keepspaces=true,
    backgroundcolor=\color{gray!5},
}
\lstset{style=mystyle}
% \geometry{margin=0in}

%%% Start Klim Packages
\usepackage{siunitx}
%%% End Klim packages

\setlist[itemize]{noitemsep, topsep=0.1pt, parsep=0.1em, leftmargin=1em}

\newcommand{\cmark}{\ding{51}} % Check mark
\newcommand{\xmark}{\ding{55}} % Cross mark
\newcommand{\approxmark}{\textbf{\LARGE \textasciitilde}}
\newcommand{\ntransforms}{28}
\newcommand{\generatorName}{COGITAO}


% ICLR TITLE 
\title{\generatorName: A Procedural Object-Centric Framework to Evaluate Compositional Generalization}
%

% NIPS TITLE
% \title{\generatorName: \\A Visual Reasoning Framework To Study
% Compositionality \& Generalization}
%

% The \author macro works with any number of authors. There are two commands
% used to separate the names and addresses of multiple authors: \And and \AND.
%
% Using \And between authors leaves it to LaTeX to determine where to break the
% lines. Using \AND forces a line break at that point. So, if LaTeX puts 3 of 4
% authors names on the first line, and the last on the second line, try using
% \AND instead of \And before the third author name.


\author{Yassine Taoudi-Benchekroun \thanks{Equal contribution. Correspondance to \href{mailto:ytaoudi@ethz.ch}{ytaoudi@ethz.ch}}\\
  Institute of Neuroinformatics \\
  ETH Zurich and University of Zurich \\
  \And
  Klim Troyan \footnotemark[1] \\
  ETH Zurich \\
  \And
  Pascal Sager \\
  Centre for Artificial Intelligence \\
  Zurich University of Applied Sciences
  \AND
  Stefan Gerber \\
  ETH Zurich
  \And
  Lukas Tuggener \\
  RWAI AG \\
  \And
  Thilo Stadelmann \\
  Centre for Artificial Intelligence \\
  Zurich University of Applied Sciences
  \AND
  Benjamin Grewe \\
  Institute of Neuroinformatics \\
  ETH Zurich and University of Zurich \\
}

\begin{document}
\maketitle

\begin{abstract}
\subfile{abstract_ICLR.tex}
\end{abstract}

\subfile{1_introduction/main.tex}
\subfile{2_related_works/main.tex}
\subfile{3_generator/main.tex}
\subfile{4_experiments/main.tex}
\subfile{5_discussion/main.tex}

\section{Reproducibility Statement}

We have made every effort to ensure that our results are fully reproducible. Detailed descriptions of the COGITAO generator, object creation process, and transformation suite are provided in Sections 3–4 of the main paper, with additional implementation details, algorithmic pseudo-code, and full experimental settings in Appendix 7–9. All training hyper-parameters, model architectures, and evaluation procedures are explicitly specified (Sections 4.1–4.4 and Appendix 8–9). To facilitate independent verification, we release the complete source code and data generation framework as anonymous supplementary material, with the same version available at our GitHub  repository \footnote{Code is available at the following URL: https://anonymous.4open.science/r/COGITAO-4E72}. The benchmark datasets used in our experiments can be reproduced or directly accessed via HuggingFace. These resources include scripts to regenerate all figures, tables, and experimental results reported in the paper. 

% \section{Acknowledgments}

% This work is part of the ‘‘Learn to learn safely’’ project funded
% by a grant of the Hasler foundation (grant nr: 21039).
% We thank Joonsu Gha for valuable guidance and feedback on the experimental section of this paper. We also thank Kevin Lopez Andrade and Michael Hodel for their insightful ideas and contributions to early versions of this work. We thank Maryam Rahimi, who reached out to us after initial release of the arxiv version, and gave us useful feedback and recommendations. We thank Matteo Saptonatti for help with the manuscript.  We are grateful to Rachida Assoughay for fruitful discussions that motivated this work. Finally, we thank Anas Ben Mejdoub for help with the visuals, and Chiara Louisa Hempel for help with framing and structuring. 


\bibliographystyle{unsrt}  
\bibliography{references}  

\clearpage
\appendix

\subfile{6_appendix/main.tex}
% \subfile{models.tex}
\clearpage
% \subfile{checklist.tex}

\end{document}